\documentclass[11pt, a4paper]{article}

% Gemeinsame Style-Definitionen (TEXINPUTS muss templates/ enthalten — siehe scripts/build_tex.py)
% bfw-style.tex — gemeinsame Style-Definitionen für alle Handouts/Aufgabenhefte
% Voraussetzung: Kompilation mit xelatex (wegen fontspec)

\usepackage[ngerman]{babel}
\usepackage{geometry}
\geometry{a4paper, top=2.2cm, bottom=2.2cm, left=2.3cm, right=2.3cm}

\usepackage{fontspec}
% Fallback-Kette: Source Sans Pro -> Calibri -> Segoe UI -> Latin Modern Sans
% (Latin Modern ist immer mit MiKTeX vorhanden)
\IfFontExistsTF{Source Sans Pro}{%
  \setmainfont{Source Sans Pro}[Ligatures=TeX]%
}{\IfFontExistsTF{Calibri}{%
  \setmainfont{Calibri}[Ligatures=TeX]%
}{\IfFontExistsTF{Segoe UI}{%
  \setmainfont{Segoe UI}[Ligatures=TeX]%
}{%
  \setmainfont{Latin Modern Sans}[Ligatures=TeX]%
}}}
\IfFontExistsTF{Source Code Pro}{%
  \setmonofont{Source Code Pro}[Scale=0.92]%
}{\IfFontExistsTF{Consolas}{%
  \setmonofont{Consolas}[Scale=0.92]%
}{%
  \setmonofont{Latin Modern Mono}[Scale=0.92]%
}}

\usepackage{xcolor}
% Farbpalette: hell, freundlich, modern
\definecolor{bfwPrimary}{HTML}{0EA5A4}    % Petrol/Türkis - Primär
\definecolor{bfwPrimaryDark}{HTML}{0F766E} % dunkler Petrol für Headings
\definecolor{bfwAccent}{HTML}{F59E0B}     % Amber - Akzent
\definecolor{bfwSuccess}{HTML}{10B981}    % Grün - Erfolg / Lösung
\definecolor{bfwWarn}{HTML}{F97316}       % Orange - Achtung
\definecolor{bfwInk}{HTML}{0F172A}        % fast Schwarz
\definecolor{bfwInkSoft}{HTML}{475569}    % gedämpftes Grau
\definecolor{bfwBgSoft}{HTML}{F8FAFC}     % off-white
\definecolor{bfwBgTeal}{HTML}{ECFEFF}     % helles Türkis
\definecolor{bfwBgAmber}{HTML}{FEF3C7}    % helles Gelb
\definecolor{bfwBgRose}{HTML}{FFE4E6}     % helles Rosé für Eselsbrücken

\usepackage[colorlinks=true, linkcolor=bfwPrimaryDark, urlcolor=bfwPrimaryDark]{hyperref}
\usepackage{enumitem}
\setlist[itemize]{leftmargin=*, itemsep=2pt, topsep=2pt}
\setlist[enumerate]{leftmargin=*, itemsep=2pt, topsep=2pt}

\usepackage[most]{tcolorbox}
\usepackage{booktabs}
\usepackage{tikz}
\usetikzlibrary{decorations.pathmorphing, positioning, shapes.geometric, arrows.meta}

% Merkkästchen — wichtige Definition / Kernpunkt
\newtcolorbox{merksatz}[1][]{%
  enhanced, breakable,
  colback=bfwBgTeal,
  colframe=bfwPrimary,
  boxrule=0.6pt,
  arc=4pt,
  left=10pt, right=10pt, top=8pt, bottom=8pt,
  fonttitle=\bfseries\color{bfwPrimaryDark},
  title={\faLightbulb\ Merksatz},
  #1
}

% Eselsbrücke — Mnemoniks
\newtcolorbox{eselsbruecke}[1][]{%
  enhanced, breakable,
  colback=bfwBgRose,
  colframe=bfwAccent,
  boxrule=0.6pt,
  arc=4pt,
  left=10pt, right=10pt, top=8pt, bottom=8pt,
  fonttitle=\bfseries\color{bfwWarn},
  title={Eselsbrücke},
  #1
}

% Beispiel-Box
\newtcolorbox{beispiel}[1][]{%
  enhanced, breakable,
  colback=bfwBgSoft,
  colframe=bfwInkSoft,
  boxrule=0.4pt,
  arc=3pt,
  left=10pt, right=10pt, top=8pt, bottom=8pt,
  fonttitle=\bfseries\color{bfwInk},
  title={Beispiel},
  #1
}

% Lösungs-Box (für Aufgabenhefte)
\newtcolorbox{loesung}[1][]{%
  enhanced, breakable,
  colback=bfwBgAmber!40,
  colframe=bfwSuccess,
  boxrule=0.6pt,
  arc=3pt,
  left=10pt, right=10pt, top=8pt, bottom=8pt,
  fonttitle=\bfseries\color{bfwSuccess!70!black},
  title={Lösung},
  #1
}

% Glossar-Eintrag
\newcommand{\glossareintrag}[2]{%
  \par\noindent\textbf{\color{bfwPrimaryDark}#1} \enspace #2\par\smallskip
}

% Section-Stil — etwas farbiger als Standard
\usepackage{titlesec}
\titleformat{\section}
  {\Large\bfseries\color{bfwPrimaryDark}}
  {\thesection}{0.6em}{}
\titleformat{\subsection}
  {\large\bfseries\color{bfwPrimary}}
  {\thesubsection}{0.5em}{}
\titleformat{\subsubsection}
  {\normalsize\bfseries\color{bfwInk}}
  {\thesubsubsection}{0.4em}{}

% TikZ-Stil "skizzenhaft" — etwas weicher als Rough.js, aber stabil
\tikzset{
  roughbox/.style = {
    draw=bfwPrimaryDark, thick, fill=bfwBgTeal,
    rounded corners=4pt, inner sep=6pt
  },
  roughaccent/.style = {
    draw=bfwAccent, thick, fill=bfwBgAmber,
    rounded corners=4pt, inner sep=6pt
  },
  roughline/.style = {
    draw=bfwInkSoft, thick, -{Stealth[length=2.5mm]}
  }
}

% Bullet-Symbole (bewusst kein FontAwesome — vermeidet Font-Installations-Probleme)
\newcommand{\faLightbulb}{$\bullet$}
\newcommand{\faBookmark}{$\bullet$}

% Header/Footer
\usepackage{fancyhdr}
\pagestyle{fancy}
\fancyhf{}
\renewcommand{\headrulewidth}{0pt}
\fancyfoot[L]{\small\color{bfwInkSoft}\thepage}
\fancyfoot[R]{\small\color{bfwInkSoft} BFW M-064 Beschaffung im Büromanagement}


% Footer für Lernfeld 1
\fancyfoot[L]{\small\color{bfwInkSoft}\copyright\ Can Siebert\,\textperiodcentered\,2026\,\textperiodcentered\,Lernfeld 1 --- Die eigene Rolle im Betrieb mitgestalten}

\title{\vspace*{-1.5cm}\Huge\bfseries\color{bfwPrimaryDark}Session~4: Den Betrieb präsentieren \& Wiederholung\\[0.4em]
  \large\color{bfwInkSoft}\textnormal{Überzeugend vortragen und Lernfeld~1 sichern}}
\author{\textsf{Lernfeld 1 --- Die eigene Rolle im Betrieb mitgestalten \textperiodcentered{} \small\copyright\ Can Siebert\,\textperiodcentered\,2026}}
\date{}

\begin{document}
\maketitle
\thispagestyle{empty}

\vspace{1em}
\begin{merksatz}[title={\bfseries Worum geht es in dieser Session?}]
Am Ende von Lernfeld~1 steht eine praktische Fähigkeit, die Sie im Beruf immer wieder
brauchen: den eigenen Betrieb überzeugend präsentieren. Wer eine Präsentation gut
vorbereitet, die richtige Zielgruppe im Blick hat, passende Medien wählt und frei
vorträgt, wirkt sicher und professionell. In diesem Handout lernen Sie, wie Sie eine
Präsentation planen, aufbauen, visualisieren und halten --- und wie Sie einen Betrieb
mit einem Unternehmenssteckbrief vorstellen. Anschließend fassen wir Lernfeld~1 noch
einmal kompakt zusammen (Sessions 1--3), damit Sie gut für die IHK-Prüfung gerüstet sind.
\end{merksatz}

\tableofcontents
\vspace{1em}
\hrule
\vspace{1em}

\section{Zweck und Ziele einer Präsentation}

Eine \textbf{Präsentation} ist die anschauliche, verständliche und überzeugende
Darstellung von Informationen vor einem Publikum. Sie unterscheidet sich vom reinen
Vortrag dadurch, dass sie \emph{Visualisierung} (Folien, Flipchart, Objekte) und die
gezielte Ansprache eines Publikums bewusst einsetzt. Typische Ziele sind:

\begin{itemize}
  \item \textbf{Informieren} --- Wissen und Fakten verständlich vermitteln.
  \item \textbf{Überzeugen} --- das Publikum für eine Idee, ein Produkt oder eine Entscheidung gewinnen.
  \item \textbf{Aktivieren} --- zum Handeln, Nachdenken oder zu einer Entscheidung anregen.
\end{itemize}

\begin{merksatz}
Eine Präsentation ist immer auf ein \emph{Ziel} und ein \emph{Publikum} ausgerichtet.
Wer weiß, \emph{wen} er \emph{wovon} überzeugen will, trifft alle weiteren Entscheidungen
(Inhalt, Sprache, Medien) leichter und richtiger.
\end{merksatz}

\section{Zielgruppen- und Adressatenanalyse}

Bevor Sie eine Präsentation vorbereiten, klären Sie, \emph{vor wem} Sie sprechen. Von der
\textbf{Zielgruppe} (auch: Adressaten) hängen Inhalt, Sprachniveau, Beispiele, Umfang und
Medienwahl ab. Leitfragen:

\begin{itemize}
  \item Wer sitzt im Publikum? (Vorwissen, Funktion, Alter, Erwartungen)
  \item Was interessiert die Gruppe, welchen Nutzen hat sie vom Thema?
  \item Wie groß ist das Publikum und wie viel Zeit steht zur Verfügung?
  \item Welche Vorbehalte oder Fragen könnten auftreten?
\end{itemize}

\begin{beispiel}
Dasselbe Thema „Unser neues Ablagesystem” präsentieren Sie vor Auszubildenden anders als
vor der Geschäftsleitung: Bei Azubis erklären Sie mehr Grundlagen mit alltagsnahen
Beispielen, bei der Geschäftsleitung betonen Sie Nutzen, Kosten und Zahlen --- kürzer und
ergebnisorientiert.
\end{beispiel}

\section{Vorbereitung: der Weg zur Präsentation}

Eine gelungene Präsentation entsteht nicht spontan, sondern durch systematische
Vorbereitung in vier Schritten:

\begin{enumerate}
  \item \textbf{Thema eingrenzen} --- Kernbotschaft festlegen, unwichtige Nebenaspekte weglassen. Weniger ist mehr.
  \item \textbf{Roten Faden entwickeln} --- eine logische, nachvollziehbare Gliederung, die alle Teile verbindet.
  \item \textbf{Informationen beschaffen} --- verlässliche Quellen nutzen, Zahlen und Beispiele sammeln, Wichtiges auswählen.
  \item \textbf{Zeit planen} --- den Ablauf auf die verfügbare Redezeit abstimmen und einüben.
\end{enumerate}

\begin{eselsbruecke}
\textbf{„Thema --- Faden --- Infos --- Zeit”} --- diese vier Schritte in genau dieser
Reihenfolge führen von der leeren Seite zur fertigen Präsentation.
\end{eselsbruecke}

\section{Aufbau: Einleitung --- Hauptteil --- Schluss}

Jede Präsentation folgt dem bewährten Dreischritt:

\begin{center}
\begin{tabular}{p{3cm} p{10.5cm}}
\toprule
\textbf{Teil} & \textbf{Aufgabe}\\
\midrule
Einleitung & Begrüßen, Thema und Ziel nennen, Interesse wecken (Frage, Zahl, Beispiel), Ablauf ankündigen.\\
Hauptteil & Kernaussagen strukturiert darstellen, mit Beispielen und Visualisierung belegen, roten Faden halten.\\
Schluss & Kernbotschaft zusammenfassen, Fazit/Appell geben, danken und Raum für Fragen bieten.\\
\bottomrule
\end{tabular}
\end{center}

\medskip
Als Richtwert entfallen Einleitung und Schluss zusammen auf rund 20--30~\% der Zeit, der
Hauptteil auf den Rest. Der \emph{erste} und der \emph{letzte} Eindruck bleiben besonders
haften --- planen Sie beide bewusst.

\section{Präsentationsmedien und Visualisierung}

Das passende Medium hängt von Publikumsgröße, Ziel und Situation ab. Kein Medium ist
universell „das beste”.

\begin{center}
\begin{tabular}{p{3.2cm} p{5cm} p{5cm}}
\toprule
\textbf{Medium} & \textbf{Vorteile} & \textbf{Nachteile}\\
\midrule
Beamer/PowerPoint & Große Gruppen, Bilder, Animationen, professioneller Eindruck & Technikabhängig, verleitet zu Textwüsten\\
Flipchart & Flexibel, live entwickelbar, kein Strom nötig & Nur kleine Gruppen, begrenzter Platz\\
Whiteboard & Spontan beschreib- und löschbar & Nicht speicherbar, geringe Reichweite\\
Pinnwand/Karten & Ideen sammeln und ordnen, alle beteiligen & Vorbereitungsintensiv, wenig Text je Karte\\
Handout & Bleibt beim Publikum, Notizen möglich, entlastet Folien & Lenkt ab, wenn zu früh verteilt\\
\bottomrule
\end{tabular}
\end{center}

\begin{merksatz}
Ein \textbf{Handout} sichert die Kerninhalte zum Nachlesen und entlastet die Folien von
Details. Verteilen Sie es in der Regel erst \emph{am Ende}, damit es während des Vortrags
nicht ablenkt.
\end{merksatz}

\section{Foliengestaltung und Gestaltungsregeln}

Folien sollen den Vortrag \emph{unterstützen}, nicht ersetzen. Bewährte Regeln:

\begin{itemize}
  \item \textbf{KISS} --- „Keep it short and simple”: so einfach und reduziert wie möglich.
  \item \textbf{6\texttimes6-Regel} --- höchstens ca.\ sechs Stichpunkte je Folie und ca.\ sechs Wörter je Stichpunkt.
  \item \textbf{Schriftgröße} --- groß und gut lesbar (Richtwert mind.\ ca.\ 24~pt).
  \item \textbf{Kontrast} --- dunkle Schrift auf hellem Grund (oder umgekehrt); keine grellen Farbkombinationen.
  \item \textbf{Einheitlichkeit} --- gleiche Schriftart, Farben und Aufbau auf allen Folien.
  \item \textbf{Bilder statt Text} --- Zahlen als Diagramm, Zusammenhänge als Schaubild.
\end{itemize}

\begin{beispiel}
Statt der Textfolie „Unser Unternehmen wurde 1998 gegründet, beschäftigt über 250
Mitarbeitende an drei Standorten und produziert hochwertige Büromöbel für Geschäftskunden”
schreiben Sie kurze Stichpunkte: „Gegründet 1998” --- „250+ Mitarbeitende” --- „3 Standorte”
--- „Büromöbel für Geschäftskunden”. Den Rest erzählen \emph{Sie}.
\end{beispiel}

\section{Vortragstechnik}

Der beste Inhalt wirkt nur, wenn der Vortrag überzeugt. Wichtige Elemente:

\begin{itemize}
  \item \textbf{Frei sprechen} --- mit Stichwortkarten statt Ablesen; wirkt lebendig und glaubwürdig.
  \item \textbf{Stimme und Tempo} --- deutlich, nicht zu schnell; Pausen bewusst einsetzen.
  \item \textbf{Blickkontakt} --- das Publikum anschauen, nicht die Folien oder den Boden.
  \item \textbf{Körpersprache} --- aufrechte, offene Haltung, ruhige Gestik, sicherer Stand.
  \item \textbf{Fragen und Feedback} --- Fragen zulassen, ruhig beantworten, Rückmeldungen annehmen.
\end{itemize}

\begin{merksatz}
\textbf{Lampenfieber} ist normal und sogar nützlich --- es macht wach. Gegenmittel: gute
Vorbereitung und Proben, tief durchatmen vor dem Start, langsam beginnen, einen freundlichen
Blickkontaktpunkt suchen und Pausen aushalten.
\end{merksatz}

\section{Einen Betrieb präsentieren: der Unternehmenssteckbrief}

Um einen Betrieb vorzustellen, nutzt man einen \textbf{Unternehmenssteckbrief} --- eine
kompakte Übersicht der wichtigsten Merkmale:

\begin{itemize}
  \item \textbf{Branche} --- in welchem Wirtschaftszweig ist das Unternehmen tätig?
  \item \textbf{Rechtsform} --- z.\,B.\ Einzelunternehmen, OHG, KG, GmbH, AG (siehe Wiederholung Session~3).
  \item \textbf{Produkte/Dienstleistungen} --- was wird angeboten?
  \item \textbf{Standort(e)} --- Sitz und weitere Niederlassungen.
  \item \textbf{Leitbild/Ziele} --- wofür steht das Unternehmen (wirtschaftliche, ökologische, soziale Ziele)?
\end{itemize}

\begin{beispiel}
\emph{Duisdorfer BüroKonzept KG}: Branche Büromöbel/Büroausstattung; Rechtsform KG;
Produkte hochwertige Büromöbel und Serviceleistungen; Standort Bonn-Duisdorf; Leitbild
u.\,a.\ Kundenservice, Umweltschutz (Recycling, Fotovoltaik) und sichere Arbeitsplätze.
\end{beispiel}

\section{Bewertung und Feedbackregeln}

Nach einer Präsentation folgt oft ein Feedback. Damit es fair und hilfreich ist, gelten
Regeln:

\begin{itemize}
  \item Erst das \textbf{Positive}, dann Verbesserungsvorschläge.
  \item \textbf{Konkret} und \textbf{beschreibend} bleiben („Die Folie war sehr textlastig”), nicht abwertend.
  \item In der \textbf{Ich-Form} sprechen („Ich hätte mir mehr Beispiele gewünscht”).
  \item \textbf{Wertschätzend} und \textbf{umsetzbar} formulieren.
  \item Wer Feedback erhält, \textbf{hört zu}, ohne sich sofort zu rechtfertigen.
\end{itemize}

Bewertungskriterien sind typischerweise: Aufbau/roter Faden, Inhalt/Verständlichkeit,
Foliengestaltung, Vortragsweise (Stimme, Blickkontakt, Körpersprache) und Zeiteinhaltung.

\section{Wiederholung: Lernfeld~1 im Überblick}

Diese Session schließt Lernfeld~1 ab. Zur Auffrischung die Kerninhalte der Sessions 1--3.

\subsection{Session~1 --- Rechtliche Grundlagen der Berufsausbildung}

Die Berufsausbildung ist im \textbf{Berufsbildungsgesetz (BBiG)} geregelt (duales System:
Betrieb + Berufsschule). Der \emph{Berufsausbildungsvertrag} begründet Rechte und
Pflichten: Auszubildende müssen lernen, sorgfältig arbeiten, Weisungen befolgen und die
Berufsschule besuchen; Ausbildende müssen die Inhalte vermitteln, Ausbildungsmittel
bereitstellen und Vergütung zahlen. Junge Beschäftigte schützt das
\textbf{Jugendarbeitsschutzgesetz (JArbSchG)}. Die \textbf{Mitbestimmung} erfolgt durch
\emph{Betriebsrat} und \emph{Jugend- und Auszubildendenvertretung (JAV)}.

\subsection{Session~2 --- Betrieb in der Wirtschaft}

Am Anfang steht das \textbf{Bedürfnis} (empfundener Mangel). Mit Kaufkraft wird daraus
\textbf{Bedarf}, am Markt geäußert \textbf{Nachfrage}. Der \textbf{einfache
Wirtschaftskreislauf} verbindet \emph{Haushalte} und \emph{Unternehmen}: Haushalte stellen
Arbeit bereit und erhalten Einkommen; Unternehmen bieten Güter und erhalten dafür Geld.
Ein Betrieb verfolgt \emph{ökonomische} (Gewinn, Produktivität, Rentabilität, Liquidität),
\emph{ökologische} und \emph{soziale} Ziele --- die in \emph{Zielharmonie} oder
\emph{Zielkonflikt} stehen können.

\subsection{Session~3 --- Unternehmen, Firma, Rechtsformen und Organisation}

Die \textbf{Firma} ist der Name, unter dem ein Kaufmann seine Geschäfte betreibt. Bei den
\textbf{Rechtsformen} unterscheidet man vor allem \emph{Personengesellschaften} (z.\,B.\
Einzelunternehmen, OHG, KG --- persönliche, oft unbeschränkte Haftung) und
\emph{Kapitalgesellschaften} (z.\,B.\ GmbH, AG --- Haftung grundsätzlich auf das
Gesellschaftsvermögen beschränkt). Die \textbf{Aufbauorganisation} (Stellen, Abteilungen,
Organigramm) legt fest, \emph{wer wofür zuständig} ist.

\begin{eselsbruecke}
\textbf{Haftung merken:} \emph{Person} haftet mit der \emph{Person} (Privatvermögen) ---
\emph{Kapital}gesellschaft haftet nur mit dem \emph{Kapital} (Gesellschaftsvermögen).
\end{eselsbruecke}

\section{Zusammenfassung}

\begin{enumerate}
  \item Eine \textbf{Präsentation} ist auf \textbf{Ziel und Zielgruppe} ausgerichtet --- die Adressatenanalyse steht am Anfang.
  \item Vorbereitung in vier Schritten: \textbf{Thema eingrenzen, roter Faden, Informationen, Zeitplanung.}
  \item Aufbau nach dem Dreischritt \textbf{Einleitung --- Hauptteil --- Schluss.}
  \item \textbf{Medien} situativ wählen (Beamer, Flipchart, Whiteboard, Pinnwand, Handout) --- jedes hat Vor- und Nachteile.
  \item Folien nach \textbf{KISS} und \textbf{6\texttimes6-Regel}: wenig Text, große Schrift, hoher Kontrast, Diagramme.
  \item \textbf{Vortragstechnik:} frei sprechen, Blickkontakt, Körpersprache, Lampenfieber beherrschen.
  \item Einen Betrieb stellt man mit einem \textbf{Unternehmenssteckbrief} vor (Branche, Rechtsform, Produkte, Standort, Leitbild).
  \item \textbf{Feedback} fair und wertschätzend geben; Lernfeld~1 (Sessions 1--3) für die IHK-Prüfung wiederholen.
\end{enumerate}

\newpage

\section*{Glossar}
\addcontentsline{toc}{section}{Glossar}

\glossareintrag{Präsentation}{Anschauliche, verständliche und überzeugende Darstellung von Informationen vor einem Publikum unter Einsatz von Visualisierung.}

\glossareintrag{Zielgruppe (Adressaten)}{Das Publikum, an das sich die Präsentation richtet; bestimmt Inhalt, Sprache, Umfang und Medienwahl.}

\glossareintrag{Adressatenanalyse}{Untersuchung von Vorwissen, Interessen, Größe und Erwartungen des Publikums vor der Vorbereitung.}

\glossareintrag{Roter Faden}{Logische, nachvollziehbare Gliederung, die alle Teile einer Präsentation sinnvoll verbindet.}

\glossareintrag{Einleitung}{Erster Teil: begrüßen, Thema und Ziel nennen, Interesse wecken, Ablauf ankündigen.}

\glossareintrag{Hauptteil}{Kernteil: die Kernaussagen strukturiert, belegt und visualisiert darstellen.}

\glossareintrag{Schluss}{Letzter Teil: Kernbotschaft zusammenfassen, Fazit/Appell, danken, Fragen zulassen.}

\glossareintrag{Beamer/PowerPoint}{Projektionsmedium für große Gruppen; ermöglicht Bilder und Animationen, ist aber technikabhängig.}

\glossareintrag{Flipchart}{Papierblock auf Ständer; flexibel und live beschreibbar, geeignet für kleine Gruppen.}

\glossareintrag{Pinnwand}{Stellwand für Moderationskarten; zum Sammeln und Ordnen von Ideen, bezieht alle ein.}

\glossareintrag{Moderationskarten}{Beschriftbare Karten für Pinnwände; kurze Aussagen, flexibel umsortierbar.}

\glossareintrag{Handout}{Begleitunterlage für das Publikum zum Nachlesen; entlastet Folien, meist erst am Ende verteilt.}

\glossareintrag{KISS-Prinzip}{„Keep it short and simple” --- Folien so einfach und reduziert wie möglich gestalten.}

\glossareintrag{6\texttimes6-Regel}{Faustregel: höchstens ca.\ sechs Stichpunkte je Folie und ca.\ sechs Wörter je Stichpunkt.}

\glossareintrag{Visualisierung}{Anschauliche Darstellung von Inhalten durch Bilder, Diagramme, Schaubilder oder Objekte.}

\glossareintrag{Lampenfieber}{Nervosität vor/während eines Auftritts; durch Vorbereitung, Atmung und Routine beherrschbar.}

\glossareintrag{Körpersprache}{Wirkung durch Haltung, Gestik, Mimik und Blickkontakt; unterstützt oder untergräbt den Vortrag.}

\glossareintrag{Unternehmenssteckbrief}{Kompakte Übersicht eines Betriebs: Branche, Rechtsform, Produkte, Standort, Leitbild/Ziele.}

\glossareintrag{Feedbackregeln}{Vereinbarungen für faire Rückmeldung: erst Positives, konkret, in Ich-Form, wertschätzend, umsetzbar.}

\end{document}
